\chapter{کاوش در اجزای سیستم}

اکنون که بر اصول پیمایش در سیستم فایل مسلط شده‌اید، زمان آن فرا رسیده است که یک کاوش هدایت‌شده را در یک توزیع مدرن لینوکس آغاز کنید. پیش از ورود به جزئیات، با سه دستور استراتژیک آشنا می‌شویم که ابزارهای کلیدی ما در این کاوش فنی خواهند بود:

\begin{itemize}
  \item \cmd{ls}: جهت فهرست کردن محتویات دایرکتوری‌ها.
  \item \cmd{file}: جهت شناسایی نوع و ماهیت فایل‌ها.
  \item \cmd{less}: جهت مشاهده محتوای فایل‌های متنی.
\end{itemize}

\section{کاربردها و گزینه‌های دستور ls}

دستور \cmd{ls} (مخفف \en{List}) یکی از پرکاربردترین فرامین در محیط خط فرمان لینوکس است؛ چراکه توانایی بالایی در استخراج اطلاعات حیاتی درباره فایل‌ها و دایرکتوری‌ها دارد. با اجرای ساده‌ی این دستور بدون هیچ آرگومانی، فهرست محتویات دایرکتوری کاری فعلی نمایش داده می‌شود که شامل تمامی فایل‌ها و دایرکتوری‌های فرعی موجود (به استثنای فایل‌های پنهان) است:

\begin{latin}
\begin{lstlisting}
[me@linuxbox ~]$ ls
Desktop  Documents  Music  Pictures  Public  Templates  Videos
\end{lstlisting}
\end{latin}

علاوه بر پایش مسیر فعلی، می‌توانید با افزودن یک مسیر به‌عنوان آرگومان به دستور \cmd{ls}، محتویات دایرکتوری دیگری را مشاهده نمایید:

\begin{latin}
\begin{lstlisting}
[me@linuxbox ~]$ ls /usr
bin  games  include  lib  local  sbin  share  src
\end{lstlisting}
\end{latin}

همچنین، دستور \cmd{ls} امکان فهرست کردن هم‌زمان چندین دایرکتوری را فراهم می‌سازد. در نمونه‌ی زیر، محتویات دایرکتوری خانگی کاربر (که با نماد \cmd{\textasciitilde} نشان داده می‌شود) و دایرکتوری \file{/usr/} به‌صورت یکپارچه نمایش داده شده‌اند:

\begin{latin}
\begin{lstlisting}
[me@linuxbox ~]$ ls ~ /usr
/home/me:
Desktop  Documents  Music  Pictures  Public  Templates  Videos

/usr:
bin  games  include  lib  local  sbin  share  src
\end{lstlisting}
\end{latin}

افزون بر این، با بهره‌گیری از گزینه‌ها (\en{Options})، می‌توان ساختار خروجی دستور \cmd{ls} را جهت نمایش جزئیات فنی بیشتر تغییر داد:

\begin{latin}
\begin{lstlisting}
[me@linuxbox ~]$ ls -l
total 56
drwxrwxr-x 2 me me 4096 2017-10-26 17:20 Desktop
drwxrwxr-x 2 me me 4096 2017-10-26 17:20 Documents
drwxrwxr-x 2 me me 4096 2017-10-26 17:20 Music
drwxrwxr-x 2 me me 4096 2017-10-26 17:20 Pictures
drwxrwxr-x 2 me me 4096 2017-10-26 17:20 Public
drwxrwxr-x 2 me me 4096 2017-10-26 17:20 Templates
drwxrwxr-x 2 me me 4096 2017-10-26 17:20 Videos
\end{lstlisting}
\end{latin}

به‌عنوان مثال، با الحاق گزینه‌ی \cmd{-l} به دستور، خروجی به حالت قالب بلند (\en{Long Format}) تغییر می‌یابد. این قالب، اطلاعات جامع و تفصیلی نظیر مجوزهای دسترسی (\en{Permissions})، مالک فایل، گروه کاربری، حجم فایل بر حسب بایت و زمان آخرین تغییرات را در اختیار کاربر قرار می‌دهد.

\section{گزینه‌ها و آرگومان‌ها (Options and Arguments)}

در این بخش به بررسی یکی از مفاهیم بنیادین در ساختار دستورات لینوکس می‌پردازیم. غالب فرامین از یک الگوی استاندارد پیروی می‌کنند که از دو مؤلفه‌ی اصلی تشکیل شده است:

\begin{enumerate}
  \item \textbf{گزینه‌ها (\en{Options}):} پارامترهایی هستند که رفتار پیش‌فرض دستور را تغییر می‌دهند. این موارد معمولاً با یک خط تیره (\cmd{-}) یا دو خط تیره (\cmd{--}) متمایز می‌شوند. برای نمونه، در دستور \cmd{ls -l}، گزینه‌ی \cmd{-l} موجب تغییر خروجی به حالت «قالب بلند» (\en{Long Format}) می‌گردد.

  \item \textbf{آرگومان‌ها (\en{Arguments}):} ورودی‌هایی هستند که هدف عملیاتی دستور را مشخص می‌کنند. در واقع دستور بر روی این موارد اجرا می‌شود. برای مثال در دستور \cmd{ls /etc}، دایرکتوری \file{/etc} آرگومانی است که به فرمان \cmd{ls} ابلاغ می‌کند محتویات این مسیر خاص را فهرست نماید.
\end{enumerate}

به‌طور کلی، فرمت استاندارد اکثر دستورات به شرح زیر است:

\begin{latin}
\begin{lstlisting}
command -options arguments
\end{lstlisting}
\end{latin}

این ساختار منعطف به شما اجازه می‌دهد تا دستورات را برای اجرای وظایف خاص، دقیقاً مطابق با نیاز خود پیکربندی کنید. گزینه‌ها در دنیای لینوکس به دو دسته‌ی کلی تقسیم می‌شوند:

\begin{enumerate}
  \item \textbf{گزینه‌های کوتاه (\en{Short Options}):} این گزینه‌ها معمولاً از یک نویسه‌ی واحد تشکیل شده و با یک خط تیره آغاز می‌شوند (مانند \cmd{-l}).

  \item \textbf{گزینه‌های بلند (\en{Long Options}):} این گزینه‌ها که در ابزارهای پروژه‌ی \en{GNU} بسیار رایج هستند، به‌صورت یک واژه‌ی کامل بوده و با دو خط تیره آغاز می‌شوند (مانند \cmd{--long-format}). گزینه‌های بلند خوانایی دستورات را به‌ویژه در اسکریپت‌نویسی ارتقا می‌دهند.
\end{enumerate}

در بسیاری از دستورات، می‌توان گزینه‌های کوتاه را به‌صورت زنجیره‌ای ادغام کرد. به عنوان مثال، برای ترکیب دو گزینه‌ی \cmd{-l} (قالب بلند) و \cmd{-t} (مرتب‌سازی بر اساس زمان)، می‌توان به جای نگارش \cmd{ls -l -t}، از فرم کوتاه‌تر \cmd{ls -lt} استفاده کرد. این دستور، محتویات را با جزئیات کامل و بر اساس زمان آخرین ویرایش (از جدید به قدیم) نمایش می‌دهد:

\begin{latin}
\begin{lstlisting}
[me@linuxbox ~]$ ls -lt
\end{lstlisting}
\end{latin}

برای معکوس کردن ترتیب نمایش، می‌توان گزینه‌ی بلند \cmd{--reverse} را به دستور افزود. این گزینه می‌تواند به‌تنهایی یا در کنار گزینه‌های کوتاه به کار رود:

\begin{latin}
\begin{lstlisting}
[me@linuxbox ~]$ ls -lt --reverse
\end{lstlisting}
\end{latin}

این مثال به‌خوبی نشان می‌دهد که چگونه با تلفیق گزینه‌های مختلف، می‌توان کنترل دقیقی بر خروجی ابزارها داشت.

\begin{note}
\textbf{نکته:} توجه داشته باشید که گزینه‌ها نیز همانند نام فایل‌ها در لینوکس، حساس به حروف بزرگ و کوچک (\en{Case-Sensitive}) هستند. برای مثال، گزینه‌ی \cmd{-r} و \cmd{-R} در بسیاری از دستورات (مانند \cmd{ls}) عملکردهایی کاملاً متفاوت و مجزا دارند؛ لذا رعایت دقیق بزرگی و کوچکی حروف در هنگام فراخوانی گزینه‌ها الزامی است.
\end{note}

دستور \cmd{ls} یکی از منعطف‌ترین ابزارهای خط فرمان لینوکس است که به واسطه‌ی گزینه‌های متعدد، امکان استخراج پایش‌های متنوعی از سیستم‌ فایل را فراهم می‌کند. در جدول زیر، رایج‌ترین گزینه‌ها و عملکرد فنی هر یک تدوین شده است:

\begin{table}[htbp]
\centering
\caption{گزینه‌های متداول دستور \cmd{ls}}
\label{tab:ls-options}
\begin{tabularx}{\textwidth}{l l X}
\toprule
\textbf{گزینه‌ی کوتاه} & \textbf{گزینه‌ی بلند} & \textbf{شرح عملکرد} \\
\midrule
\cmd{-a} & \cmd{--all} & فهرست کردن تمامی فایل‌ها، از جمله فایل‌های پنهان (مواردی که با \cmd{.} آغاز می‌شوند). \\
\addlinespace
\cmd{-A} & \cmd{--almost-all} & مشابه گزینه \cmd{-a} با این تفاوت که از نمایش دایرکتوری جاری (\cmd{.}) و والد (\cmd{..}) صرف‌نظر می‌کند. \\
\addlinespace
\cmd{-d} & \cmd{--directory} & به‌جای نمایش محتویات یک دایرکتوری، اطلاعات مربوط به خودِ دایرکتوری را نمایش می‌دهد (غالباً همراه با \cmd{-l} استفاده می‌شود). \\
\addlinespace
\cmd{-F} & \cmd{--classify} & الحاق یک کاراکتر نمادین به انتهای نام‌ها برای تشخیص سریع ماهیت فایل (مثلاً \cmd{/} برای دایرکتوری‌ها). \\
\addlinespace
\cmd{-h} & \cmd{--human-readable} & در قالب بلند (\cmd{-l})، حجم فایل‌ها را با واحدهای قابل‌درک (مانند \en{KB}، \en{MB}، \en{GB}) نمایش می‌دهد. \\
\addlinespace
\cmd{-l} & --- & نمایش خروجی در قالب بلند (\en{Long Format}) شامل تمامی فراداده‌ها و مجوزهای امنیتی. \\
\addlinespace
\cmd{-r} & \cmd{--reverse} & معکوس کردن ترتیب نمایش (به‌صورت پیش‌فرض، چیدمان بر اساس حروف الفباست). \\
\addlinespace
\cmd{-S} & --- & مرتب‌سازی خروجی بر اساس حجم فایل‌ها (از بزرگ به کوچک). \\
\addlinespace
\cmd{-t} & --- & مرتب‌سازی بر اساس زمان آخرین ویرایش (از جدیدترین به قدیمی‌ترین). \\
\bottomrule
\end{tabularx}
\end{table}

بهره‌گیری هوشمندانه از این گزینه‌ها به شما این امکان را می‌دهد تا در کمترین زمان ممکن، دقیق‌ترین جزئیات فنی را از ساختار سیستم فایل استخراج نمایید.

\section{تحلیل قالب بلند (-l)}

همان‌طور که پیش‌تر اشاره شد، گزینه \cmd{-l} خروجی دستور \cmd{ls} را به ساختار موسوم به قالب بلند (\en{Long Format}) تغییر می‌دهد. این قالب، فراداده‌های (\en{Metadata}) حیاتی هر فایل و دایرکتوری را با جزئیات دقیق در اختیار کاربر قرار می‌دهد. برای کالبدشکافی این داده‌ها، خروجی نمونه‌ی زیر را در نظر بگیرید:

\begin{latin}
\begin{lstlisting}
-rw-r--r-- 1 root root 3576296 2017-04-03 11:05 Experience ubuntu.ogg
-rw-r--r-- 1 root root 1186219 2017-04-03 11:05 kubuntu-leaflet.png
-rw-r--r-- 1 root root   47584 2017-04-03 11:05 logo-Edubuntu.png
-rw-r--r-- 1 root root   44355 2017-04-03 11:05 logo-Kubuntu.png
-rw-r--r-- 1 root root   34391 2017-04-03 11:05 logo-Ubuntu.png
-rw-r--r-- 1 root root   32059 2017-04-03 11:05 oo-cd-cover.odf
-rw-r--r-- 1 root root  159744 2017-04-03 11:05 oo-derivatives.doc
-rw-r--r-- 1 root root   27837 2017-04-03 11:05 oo-maxwell.odt
-rw-r--r-- 1 root root   98816 2017-04-03 11:05 oo-trig.xls
-rw-r--r-- 1 root root  453764 2017-04-03 11:05 oo-welcome.odt
-rw-r--r-- 1 root root  358374 2017-04-03 11:05 ubuntu Sax.ogg
\end{lstlisting}
\end{latin}

هر ستون از این خروجی حامل اطلاعات فنی مشخصی است که در ادامه به تشریح دقیق هر یک خواهیم پرداخت.

خروجی در وضعیت قالب بلند (\en{Long Format})، شامل چندین ستون محاسباتی و توصیفی است که شناسنامه فنی هر فایل یا دایرکتوری را تشکیل می‌دهند. در جدول زیر، جزئیات دقیق هر یک از این فیلدها تشریح شده است:

\begin{table}[htbp]
\centering
\caption{تحلیل فیلدهای خروجی \cmd{ls -l}}
\label{tab:ls-l-fields}
\begin{tabularx}{\textwidth}{l X}
\toprule
\textbf{فیلد خروجی} & \textbf{مفهوم و عملکرد فنی} \\
\midrule
\cmd{-rw-r--r--} & این رشته‌ی ده کاراکتری، سطوح دسترسی (\en{Permissions}) و نوع فایل را مشخص می‌کند.
\begin{itemize}
  \item کاراکتر اول: بیانگر ماهیت فایل است؛ \cmd{-} برای فایل‌های معمولی و \cmd{d} برای دایرکتوری‌ها.
  \item نه کاراکتر بعدی: به سه بلوک سه‌تایی تقسیم می‌شوند که به ترتیب مجوزهای مالک (\en{Owner})، گروه (\en{Group}) و سایر کاربران (\en{Others}) را تعیین می‌کنند. این مجوزها شامل خواندن (\cmd{r})، نوشتن (\cmd{w}) و اجرا (\cmd{x}) هستند.
\end{itemize} \\
\addlinespace
\cmd{1} & تعداد پیوندهای سخت (\en{Hard Links}): تعداد ارجاعات مستقیم سیستم فایل به این داده. (برای درک عمیق‌تر، مبحث پیوندها را در فصول آتی دنبال کنید). \\
\addlinespace
\cmd{root} & مالک فایل (\en{Owner}): شناسه کاربری (\en{UID}) که مالکیت حقوقی فایل را در اختیار دارد. \\
\addlinespace
\cmd{root} & گروه مالک (\en{Group}): شناسه گروه (\en{GID}) که فایل به آن منتسب شده است. \\
\addlinespace
\cmd{32059} & حجم فایل (\en{Size}): اندازه واقعی فایل بر حسب بایت. \\
\addlinespace
\cmd{2017-04-03 11:05} & زمان آخرین ویرایش (\en{Mtime}): دقیق‌ترین زمان ثبت‌شده برای آخرین تغییر در محتوای فایل. \\
\addlinespace
\cmd{oo-cd-cover.odf} & نام فایل (\en{Filename}): شناسه یا نام اختصاصی فایل یا دایرکتوری. \\
\bottomrule
\end{tabularx}
\end{table}

درک این ستون‌ها به شما اجازه می‌دهد تا علاوه بر مدیریت فضا، امنیت و دسترسی‌های لایه‌بندی شده را در محیط لینوکس به‌دقت پایش کنید.

\section{شناسایی نوع فایل‌ها با دستور file}

در فرآیند پایش سیستم، تشخیص دقیق نوع داده‌های ذخیره‌شده در یک فایل اهمیت حیاتی دارد. همان‌طور که پیش‌تر اشاره شد، در لینوکس برخلاف سایر سیستم‌عامل‌ها، پسوند فایل (\en{File Extension}) لزوماً بازتاب‌دهنده‌ی ماهیت درونی آن نیست؛ برای نمونه، فایلی با نام \file{picture.jpg} ممکن است برخلاف ظاهرش، یک تصویر با فرمت \en{JPEG} نباشد.

دستور \cmd{file} ابزاری تخصصی برای رفع این ابهام است. این دستور با تحلیل ساختار درونی فایل، نوع واقعی محتوا را استخراج می‌کند:

\begin{latin}
\begin{lstlisting}
file filename
\end{lstlisting}
\end{latin}

پس از فراخوانی دستور \cmd{file} بر روی یک هدف مشخص، توصیفی فنی از ماهیت آن ارائه می‌گردد. این ابزار به جای اتکا به نام فایل، از روش‌هایی نظیر بررسی اعداد جادویی (\en{Magic Numbers}) و تحلیل هدرهای داده برای تشخیص فرمت استفاده می‌کند.

به‌عنوان مثال:

\begin{latin}
\begin{lstlisting}
[me@linuxbox ~]$ file picture.jpg
picture.jpg: JPEG image data, JFIF standard 1.01
\end{lstlisting}
\end{latin}

در سیستم‌عامل‌های شبه‌یونیکس نظیر لینوکس، اصلی بنیادین حاکم است: «همه چیز یک فایل است» (\en{Everything is a file}). این مفهوم به این معناست که در لایه‌های انتزاعی سیستم، منابع سخت‌افزاری، فرآیندهای در حال اجرا و حتی دایرکتوری‌ها، همگی به‌عنوان «فایل» مدل‌سازی شده و از طریق رابط‌های استاندارد مدیریت فایل قابل دسترسی هستند. با پیشرفت در مسیر یادگیری، ابعاد فنی و کاربردی این معماری برای شما شفاف‌تر خواهد شد.

اگرچه بسیاری از فایل‌ها دارای فرمت‌های استانداردی همچون \en{MP3} یا \en{JPEG} هستند، اما در محیط لینوکس با دسته‌بندی‌های دیگری از فایل‌ها (مانند فایل‌های دستگاه بلوکی، فایل‌های دستگاه کاراکتری و سوکت‌ها) مواجه خواهید شد که ساختاری غیرمعمول و عملکردی پیچیده دارند. درک این تنوع ساختاری، کلید تعامل حرفه‌ای با هسته‌ی سیستم‌عامل است.

\section{مشاهده و پیمایش محتوای متنی با ابزار less}

دستور \cmd{less} یک نمایش‌دهنده متنی (\en{Pager}) است که به‌منظور پایش بهینه و سریع فایل‌ها، به‌ویژه در ابعاد بسیار حجیم، توسعه یافته است. این ابزار به جای بارگذاری کامل داده‌ها در حافظه اصلی (\en{RAM})، محتوا را به‌صورت قطعه‌قطعه و صفحه‌به‌صفحه فراخوانی می‌کند که این رویکرد، پایداری و کارایی سیستم را در مواجهه با فایل‌های چند گیگابایتی تضمین می‌نماید.

در اکوسیستم لینوکس، اکثر فایل‌های پیکربندی (\en{Configuration Files})، فایل‌های گزارش (\en{Log Files}) و اسکریپت‌های سیستمی در قالب متن ساده (\en{Plain Text}) نگهداری می‌شوند. ابزار \cmd{less} بستری ایده‌آل برای بررسی و تحلیل این مستندات فراهم می‌آورد، بدون آنکه ریسک تغییر ناخواسته در محتوای آن‌ها وجود داشته باشد.

نام این دستور در واقع یک شوخی فنی با ابزار قدیمی‌تر \cmd{more} است؛ با این پیام که «\cmd{less} قابلیت‌های بیشتری نسبت به \cmd{more} دارد». این ابزار تجربه‌ای روان‌تر و امکانات پیمایش پیشرفته‌تری را در اختیار کاربر قرار می‌دهد.

\section{مفهوم متن ساده (Plain Text) در رایانش}

متن در دنیای رایانه، روشی بی‌مانند و کارآمد برای نمایش و ذخیره‌سازی اطلاعات است. با وجود اینکه تمامی داده‌های دیجیتال در لایه‌های زیرین به‌صورت اعداد دودویی (\en{Binary}) ذخیره می‌شوند، «متن» یکی از ابتدایی‌ترین الگوهای نگاشت (\en{Mapping}) مستقیم نویسه‌ها به کدهای عددی محسوب می‌شود.

در ساختار متنی، هر نویسه (شامل حروف، اعداد و نمادها) بر اساس استانداردهای رمزگذاری نظیر \en{ASCII} یا \en{UTF-8} به یک شناسه‌ی عددی منحصربه‌فرد تخصیص می‌یابد. این تناظر یک‌به‌یک موجب فشردگی حداکثری داده‌ها می‌گردد؛ به‌طوری که ۵۰ نویسه‌ی متنی دقیقاً معادل ۵۰ بایت داده در حافظه اشغال می‌کند.

باید میان متن ساده (\en{Plain Text}) و فایل‌های تولیدشده توسط واژه‌پردازها (مانند \en{Microsoft Word} یا \en{LibreOffice Writer}) تفکیک قائل شد. اسناد واژه‌پرداز علاوه بر نویسه‌های متنی، حاوی فراداده‌های حجیمی شامل کدهای قالب‌بندی (\en{Formatting})، جداول، تصاویر و ساختار سند هستند. در مقابل، فایل‌های متنی صرفاً شامل نویسه‌های استاندارد و تعداد محدودی کدهای کنترلی پایه (مانند \en{Tab} برای فاصله‌گذاری و \en{Newline} برای خط جدید) می‌باشند.

در معماری لینوکس، استفاده از فایل‌های متنی یک اصل استراتژیک و بنیادین است. بخش وسیعی از فایل‌های پیکربندی که وظیفه‌ی ذخیره‌ی تنظیمات حیاتی سیستم را بر عهده دارند، در این قالب نگهداری می‌شوند. همچنین، اکثر اسکریپت‌های مدیریتی و کدهای منبع برنامه‌ها نیز در قالب متن ساده نوشته شده‌اند.

تسلط بر مطالعه و ویرایش این فایل‌ها به شما امکان می‌دهد تا درک عمیق‌تری از لایه‌های درونی سیستم پیدا کنید و رفتار سیستم‌عامل را سفارشی‌سازی نمایید. در فصول آتی، با ابزارهای ویرایش متن جهت تغییر تنظیمات و توسعه اسکریپت‌ها آشنا خواهید شد؛ اما در این مرحله، تمرکز ما صرفاً بر پایش و مشاهده‌ی محتوای آن‌ها است.

فراخوانی دستور \cmd{less} به همراه یک یا چند نام فایل به‌عنوان آرگومان، محتوای آن‌ها را در یک محیط تعاملی جهت مطالعه نمایش می‌دهد:

\begin{latin}
\begin{lstlisting}
less filename
\end{lstlisting}
\end{latin}

پس از اجرای این دستور، رابط کاربری \cmd{less} فعال شده و امکان پیمایش (\en{Scrolling}) در سراسر فایل را فراهم می‌کند. به‌عنوان نمونه، برای بررسی محتوای فایل \file{/etc/passwd} که حاوی فراداده‌های مربوط به حساب‌های کاربری سیستم است، دستور زیر را اجرا کنید:

\begin{latin}
\begin{lstlisting}
[me@linuxbox ~]$ less /etc/passwd
\end{lstlisting}
\end{latin}

استفاده از این ابزار به شما اجازه می‌دهد تا بدون نگرانی از تغییر ناخواسته‌ی فایل‌های حساس سیستمی، ساختار و محتوای آن‌ها را تحلیل کنید.

جهت پیمایش حرفه‌ای و کارآمد در فایل‌های متنی با استفاده از ابزار \cmd{less}، مجموعه‌ای از دستورات صفحه‌کلید تعبیه شده است که امکان پیمایش سریع، جستجوی دقیق و مدیریت جریان نمایش داده‌ها را فراهم می‌کنند. در جدول زیر، پرتکرارترین دستورات عملیاتی این ابزار طبقه‌بندی شده است:

\begin{table}[htbp]
\centering
\caption{دستورات \cmd{less}}
\label{tab:less-commands}
\begin{tabularx}{\textwidth}{l X}
\toprule
\textbf{دستور} & \textbf{عملکرد فنی} \\
\midrule
\en{Page Up} یا \cmd{b} & پیمایش به عقب (صعودی) به اندازه‌ی یک صفحه‌ی کامل. \\
\addlinespace
\en{Page Down} یا \cmd{space} & پیمایش به جلو (نزولی) به اندازه‌ی یک صفحه‌ی کامل. \\
\addlinespace
کلید جهت‌نمای بالا & پیمایش به سمت ابتدای فایل به اندازه‌ی یک سطر. \\
\addlinespace
کلید جهت‌نمای پایین & پیمایش به سمت انتهای فایل به اندازه‌ی یک سطر. \\
\addlinespace
\cmd{G} & انتقال آنی مکان‌نما به آخرین خط فایل. \\
\addlinespace
\cmd{1G} یا \cmd{g} & انتقال آنی مکان‌نما به نخستین خط فایل. \\
\addlinespace
\cmd{/characters} & جستجوی رو به جلو (\en{Forward Search}) برای یافتن اولین رخداد رشته متنی مورد نظر. \\
\addlinespace
\cmd{n} & تکرار عملیات جستجو برای یافتن نمونه‌ی بعدی از عبارت هدف. \\
\addlinespace
\cmd{h} & فراخوانی راهنما (\en{Help}) جهت مشاهده تمامی دستورات موجود. \\
\addlinespace
\cmd{q} & خروج از محیط تعاملی و بازگشت به خط فرمان (\en{Quit}). \\
\bottomrule
\end{tabularx}
\end{table}

\section{تحلیل ابزار less و اصل Less is More}

ابزار \cmd{less} به‌عنوان جایگزینی تکامل‌یافته و قدرتمند برای برنامه‌ی قدیمی \cmd{more} توسعه یافته است. نام این ابزار، ارجاعی هوشمندانه به گزین‌گویه‌ی معروف «\en{Less is More}» (کمتر، بیشتر است) دارد که یکی از اصول بنیادین در مکتب طراحی مدرن محسوب می‌شود. هر دو ابزار در رده‌ی نمایش‌دهنده‌های متنی (\en{Pager}) قرار می‌گیرند؛ نرم‌افزارهایی که برای مطالعه‌ی تسهیل‌شده‌ی اسناد متنی طولانی به‌صورت صفحه‌به‌صفحه طراحی شده‌اند. برخلاف \cmd{more} که در نسخه‌های اولیه تنها امکان پیمایش یک‌طرفه (رو به جلو) را داشت، \cmd{less} قابلیت‌های فنی پیشرفته‌تری را ارائه می‌دهد که از آن جمله می‌توان به موارد زیر اشاره کرد:

\begin{itemize}
  \item ناوبری دوطرفه: آزادی عمل در پیمایش متن هم به سمت جلو و هم به سمت عقب.
  \item جست‌وجوی بهینه: برخورداری از موتور جست‌وجوی سریع و منعطف در میان محتوای متنی.
  \item بهره‌وری در مدیریت حافظه: توانایی بازگشایی و مدیریت فایل‌های حجیم بدون نیاز به بارگذاری کامل آن‌ها در حافظه‌ی \en{RAM}، که مانع از کندی سیستم می‌شود.
\end{itemize}

این ویژگی‌های استراتژیک، \cmd{less} را به ابزاری بی‌بدیل برای تحلیل دقیق گزارش‌های سیستمی و بررسی مستندات فنی تبدیل کرده است.

\section{کاوش در ساختار سلسله‌مراتبی سیستم فایل}

سازماندهی دایرکتوری‌ها و فایل‌ها در لینوکس از استانداردی به نام استاندارد سلسله‌مراتب سیستم فایل (\en{FHS}) پیروی می‌کند. این استاندارد چارچوبی دقیق برای جانمایی داده‌ها تعریف می‌کند که درک عمیق آن، کلید شناسایی عملکرد هسته و محل استقرار فایل‌های حیاتی سیستم است. شایان ذکر است که لزوماً تمامی توزیع‌های لینوکسی به‌صورت مطلق از این قواعد تبعیت نمی‌کنند، اما اکثریت قریب به اتفاق آن‌ها با این الگو همسو هستند.

در ادامه، با انجام یک پیمایش عملی در لایه‌های مختلف سیستم‌ فایل، مهارت‌های ناوبری خود را در محیط عملیاتی تقویت کنید. همان‌طور که پیش‌تر اشاره شد، بخش اعظمی از فایل‌های زیرساختی سیستم در قالب متن ساده (\en{Plain Text}) و با قابلیت خوانش انسانی ذخیره شده‌اند. این گشت‌و‌گذار فرصتی مغتنم برای پایش مستقیم این مستندات و درک سازوکار اجزای درونی لینوکس فراهم می‌آورد.

برای دستیابی به درک عمیق از ساختار سیستم فایل، در طول فرآیند کاوش، الگوریتم عملیاتی زیر را دنبال کنید:

\begin{enumerate}
  \item با بهره‌گیری از دستور \cmd{cd} در دایرکتوری‌های مختلف جابه‌جا شوید.
  \item از دستور \cmd{ls -l} برای مشاهده محتویات در قالب بلند استفاده کنید تا به فراداده‌های دقیق فایل‌ها دسترسی یابید.
  \item با استفاده از ابزار \cmd{file}، نوع فایل‌های مورد نظر را شناسایی کنید.
  \item چنانچه ماهیت فایل از نوع متنی (\en{text}) بود، با استفاده از \cmd{less} محتوای آن را مورد مطالعه قرار دهید.
  \item در صورتی که به‌اشتباه یک فایل دودویی (\en{Binary}) را باز کردید و کاراکترهای خروجی ترمینال ناخوانا شدند، با اجرای دستور \cmd{reset} وضعیت شبیه‌ساز ترمینال را به حالت استاندارد بازگردانید.
\end{enumerate}

\begin{note}
\textbf{نکته:} جهت پیشگیری از خطاهای تایپی، از قابلیت کپی و چسباندن سریع استفاده کنید؛ نام فایل را با دوبار کلیک انتخاب کرده و سپس با فشردن دکمه میانی ماوس، آن را در اعلان خط فرمان درج نمایید.
\end{note}

در طول این فرآیند کاوش، با اطمینان کامل عمل کنید. در لینوکس، کاربران عادی به‌صورت پیش‌فرض فاقد مجوزهای لازم برای اعمال تغییرات مخرب در سطح سیستم هستند. در این محیط هیچ لایه‌ی پنهانی وجود ندارد؛ سیستم برای کشف و یادگیری در اختیار شماست.

ساختار دایرکتوری‌ها در لینوکس، که عمدتاً بر اساس استاندارد سلسله‌مراتب سیستم فایل (\en{FHS}) شکل گرفته است، نقشه‌ای منطقی و منسجم برای سازماندهی داده‌ها فراهم می‌کند. این چیدمان که با دقت مهندسی شده است، به مدیران سیستم و کاربران اجازه می‌دهد تا با سرعت و اطمینان بالا به منابع حیاتی سیستم دسترسی پیدا کنند. در جدول زیر، به تشریح دقیق مهم‌ترین دایرکتوری‌های سطح بالا (\en{top-level directories}) در ساختار ریشه (\cmd{/}) و عملکرد تخصصی هر یک می‌پردازیم:

\begin{longtable}{l p{0.72\textwidth}}
\caption{ساختار و وظایف دایرکتوری‌های اصلی لینوکس}
\label{tab:linux-directories} \\
\toprule
\textbf{دایرکتوری} & \textbf{شرح عملکرد و محتوا} \\
\midrule
\endfirsthead

\multicolumn{2}{c}{\tablename\ \thetable{} (ادامه) \rule[-0.4em]{0pt}{1.4em}} \\
\toprule
\textbf{دایرکتوری} & \textbf{شرح عملکرد و محتوا} \\
\midrule
\endhead

\midrule
\multicolumn{2}{r}{\footnotesize ادامه در صفحه‌ی بعد \ldots} \\
\endfoot

\bottomrule
\endlastfoot

\cmd{/} & دایرکتوری ریشه (\en{Root Directory}): بالاترین سطح در ساختار سیستم فایل لینوکس. تمام فایل‌ها و دایرکتوری‌ها به‌صورت یک ساختار درختی از این نقطه آغاز می‌شوند. در موارد زیر، اسلش ابتدایی (\cmd{/}) یعنی آن دایرکتوری مستقیماً درون دایرکتوری ریشه (\cmd{/}) قرار دارد. این دایرکتوری‌ها، بخشی از ساختار سلسله‌مراتبی استاندارد لینوکس هستند که بر اساس \en{FHS} مهندسی شده‌اند: \\
\addlinespace
\cmd{/bin} & باینری‌های ضروری (\en{Essential Binaries}): حاوی فایل‌های اجرایی حیاتی برای بوت و عملکردهای پایه سیستم. در توزیع‌های نوین، این مسیر معمولاً یک پیوند نمادین (\en{Symlink}) به \file{/usr/bin} است. \\
\addlinespace
\cmd{/boot} & فایل‌های راه‌اندازی (\en{Boot Files}): این دایرکتوری میزبان تمامی مولفه‌های لازم برای راه‌اندازی سیستم است. شامل هسته لینوکس (\file{vmlinuz})، ایمیج رم‌دیسک آغازین (\file{initrd.img} یا \file{initramfs})، و فایل‌های پیکربندی بوت‌لودر (مانند \file{grub.cfg}). این دایرکتوری برای فرآیند بوت حیاتی است. \textbf{فایل‌های کلیدی:}
\begin{itemize}
  \item \file{/boot/grub/grub.cfg}: فایل پیکربندی بوت‌لودر \en{GRUB} که وظیفه‌ی بارگذاری هسته در حافظه را بر عهده دارد.
  \item \file{/boot/vmlinuz}: هسته (\en{Kernel}) فشرده‌شده‌ی لینوکس؛ قلب تپنده سیستم که مدیریت تمامی منابع سخت‌افزاری، زمان‌بندی پردازش‌ها، حافظه و ارتباط بین نرم‌افزار و سخت‌افزار را بر عهده دارد.
\end{itemize} \\
\addlinespace
\cmd{/dev} & گره‌های دستگاه (\en{Device Nodes}): این دایرکتوری شامل فایل‌های ویژه‌ای است که به عنوان رابط بین نرم‌افزارها و سخت‌افزار عمل می‌کنند. هر فایل در \file{/dev} نماینده‌ی یک دستگاه سخت‌افزاری (مانند هارددیسک، ترمینال، یا ماوس) یا یک درایور هسته است. این فایل‌ها در حقیقت فایل‌های فیزیکی روی دیسک نیستند؛ بلکه توسط هسته در زمان اجرا (به‌صورت مجازی) ایجاد می‌شوند. از این رو، \file{/dev} یک سیستم فایل مجازی است و محتوای آن با راه‌اندازی مجدد سیستم بازسازی می‌شود. \textbf{انواع گره‌های دستگاه:}
\begin{itemize}
  \item دستگاه‌های بلوکی (\en{Block Devices}): داده‌ها را به‌صورت بلوک‌هایی با اندازه‌ی ثابت خوانده و می‌نویسند. معمولاً برای هارددیسک‌ها و حافظه‌های \en{USB} استفاده می‌شوند. مثال: \file{/dev/sda} (نخستین هارددیسک شناسایی‌شده توسط سیستم).
  \item دستگاه‌های کاراکتری (\en{Character Devices}): داده‌ها را به‌صورت جریانی از بایت‌ها (کاراکتر به کاراکتر) منتقل می‌کنند. معمولاً برای دستگاه‌هایی مانند ترمینال‌ها، ماوس و چاپگرها استفاده می‌شوند. مثال: \file{/dev/tty} (ترمینالِ کنترل‌کننده) و \file{/dev/null} (سیاه‌چاله‌ی داده‌ها).
\end{itemize}
\textbf{نکته:} علاوه بر این دو نوع گره دستگاه، فایل‌های ویژه دیگری مانند پایپ‌های نام‌گذاری‌شده (\en{Named Pipes / FIFOs}) و سوکت‌های دامنه یونیکس (\en{Unix Domain Sockets}) نیز ممکن است در \file{/dev} وجود داشته باشند، اما آن‌ها گره سخت‌افزاری محسوب نمی‌شوند و برای ارتباط بین فرایندی (\en{IPC}) کاربرد دارند.
\textbf{اهمیت این دایرکتوری:} وجود \file{/dev} تجسم عینی اصلی بنیادین در خانوادهٔ سیستم‌عامل‌های شبه‌یونیکس است: «همه چیز یک فایل است» (\en{Everything is a file}). این بدان معناست که نرم‌افزارها می‌توانند با دستگاه‌های سخت‌افزاری دقیقاً به همان روشی تعامل کنند که با فایل‌های معمولی تعامل می‌کنند (با استفاده از فراخوان‌های سیستمی استانداردی مانند \cmd{read}، \cmd{write} و \cmd{open}). برای مثال، نوشتن داده در \file{/dev/sda} معادل نوشتن مستقیم و بدون واسطه روی بخش‌های فیزیکی هارددیسک است. \\
\addlinespace
\cmd{/etc} & پیکربندی‌های سراسری (\en{System Configuration}): مرکز اصلی تنظیمات سیستم که شامل اغلب فایل‌های پیکربندی سیستم و سرویس‌ها (به‌صورت متن ساده) است. \textbf{فایل‌های کلیدی:}
\begin{itemize}
  \item \file{/etc/crontab}: جدول زمان‌بندی وظایف خودکار (\en{Cron Jobs}).
  \item \file{/etc/fstab}: جدول نگاشت پارتیشن‌ها و نقاط اتصال (\en{Mount Points}).
  \item \file{/etc/passwd}: پایگاه‌داده‌ی متنی حاوی فراداده‌های حساب‌های کاربری (\en{UID}، \en{GID} و مسیر \en{Home}).
\end{itemize} \\
\addlinespace
\cmd{/home} & دایرکتوری‌های خانگی کاربران (\en{Home Directories}): محل استقرار دایرکتوری‌های شخصی کاربران عادی. امنیت سیستم ایجاب می‌کند که هر کاربر به‌طور پیش‌فرض تنها در شاخه‌ی اختصاصی خود مجوز نگارش داشته باشد. \\
\addlinespace
\cmd{/lib} & کتابخانه‌های مشترک (\en{Shared Libraries}): کتابخانه‌های ضروری برای باینری‌های موجود در \file{/bin} و \file{/sbin}. همچنین شامل ماژول‌های هسته (\file{/lib/modules}) و در برخی توزیع‌ها، \en{firmware} دستگاه‌ها. \\
\addlinespace
\cmd{/lost+found} & اشیاء بازیابی‌شده (\en{Recovered Files}): مختص سیستم‌های فایل خانواده \en{ext}. این دایرکتوری محل ذخیرهٔ قطعات یا فایل‌های کامل اما بی‌نامی است که ابزار \cmd{fsck} پس از یک خرابی ناگهانی سیستم، قادر به بازگرداندن آن‌ها به مکان اصلی‌شان نبوده است. \\
\addlinespace
\cmd{/media} & رسانه‌های قابل‌حمل (\en{Removable Media}): نقطه‌ی اتصال پیش‌فرض برای دستگاه‌های ذخیره‌سازی خارجی مانند حافظه‌های \en{USB}، دیسک‌های نوری و کارت‌های حافظه (معمولاً به‌صورت خودکار توسط سیستم). \\
\addlinespace
\cmd{/mnt} & اتصال دستی (\en{Manual Mount}): دایرکتوری سنتی برای اتصال دستی سیستم‌های فایل به‌صورت موقت، که معمولاً نیاز به دسترسی ابرکاربر دارد. \\
\addlinespace
\cmd{/opt} & نرم‌افزارهای جانبی (\en{Optional Software}): محل نصب بسته‌های نرم‌افزاری شخص ثالث که خارج از مخازن رسمی توزیع نصب می‌شوند (مانند برنامه‌های تجاری). هر برنامه معمولاً در یک زیرشاخه مجزا قرار می‌گیرد. \\
\addlinespace
\cmd{/proc} & سیستم فایل فرآیندها (\en{Process Filesystem}): یک رابط مجازی که اطلاعات لحظه‌ای دربارهٔ هسته، فرآیندهای در حال اجرا و سخت‌افزار را در قالب فایل‌های متنی ارائه می‌دهد (مانند \file{/proc/cpuinfo} و \file{/proc/meminfo}). \\
\addlinespace
\cmd{/root} & دایرکتوری خانگی ابرکاربر (\en{Root}): دایرکتوری اختصاصی کاربر \cmd{root} که جهت حفظ امنیت، مجزا از شاخه‌ی \file{/home} قرار دارد. \\
\addlinespace
\cmd{/run} & داده‌های زمان اجرا (\en{Runtime Data}): یک سیستم فایل موقت در حافظه (\en{RAM}) که اطلاعات مربوط به فرآیندها و سرویس‌های در حال اجرا، از زمان بوت تا خاموشی سیستم، در آن ذخیره می‌شود. \\
\addlinespace
\cmd{/sbin} & باینری‌های سیستمی (\en{System Binaries}): ابزارهای اجرایی که عمدتاً برای مدیریت سیستم و عیب‌یابی استفاده می‌شوند و معمولاً نیاز به دسترسی ابرکاربر دارند (مانند \cmd{fdisk}، \cmd{mount}، \cmd{ifconfig}). \\
\addlinespace
\cmd{/sys} & سیستم فایل هسته (\en{Sysfs}): یک رابط مجازی برای نمایش و کنترل پارامترهای سخت‌افزاری، درایورها و تنظیمات هسته در زمان اجرا (مشابه \file{/proc}، اما ساختاریافته‌تر). \\
\addlinespace
\cmd{/tmp} & فایل‌های موقت (\en{Temporary Files}): فضایی برای ذخیره‌سازی فایل‌های موقتی که توسط برنامه‌ها و کاربران ایجاد می‌شوند. معمولاً پس از هر بار راه‌اندازی سیستم، این دایرکتوری تخلیه می‌شود. \\
\addlinespace
\cmd{/usr} & منابع کاربری (\en{User System Resources}): یکی از بزرگ‌ترین و مهم‌ترین دایرکتوری‌ها که شامل برنامه‌های کاربردی، کتابخانه‌ها، مستندات و فایل‌های اشتراکی سیستم است. \\
\addlinespace
\cmd{/usr/bin} & باینری‌های کاربری (\en{User Binaries}): شامل اکثر دستورات و برنامه‌های قابل اجرا برای کاربران عادی (مانند \cmd{ls}، \cmd{cp}، \cmd{grep} و ویرایشگرها). \\
\addlinespace
\cmd{/usr/lib} & کتابخانه‌های کاربری (\en{User Libraries}): شامل کتابخانه‌های مشترک و ماژول‌های مورد نیاز برای برنامه‌های موجود در \file{/usr/bin} و \file{/usr/sbin}. \\
\addlinespace
\cmd{/usr/local} & سلسله‌مراتب محلی (\en{Local Hierarchy}): دایرکتوری مخصوص نصب نرم‌افزارهایی که خارج از مخازن رسمی توزیع و به‌صورت دستی (مثلاً از طریق کامپایل کد منبع (\en{Source Code})) در سیستم مستقر می‌شوند تا با فایل‌های اصلی سیستم تداخل نداشته باشند. ساختار آن مشابه \file{/usr} است (\file{/usr/local/bin}، \file{/usr/local/lib} و \ldots). \\
\addlinespace
\cmd{/usr/sbin} & باینری‌های سیستمی ثانویه (\en{System Binaries}): ابزارهای مدیریتی غیرحیاتی که برای کارهای پیشرفته‌تر سیستم استفاده می‌شوند و معمولاً نیاز به دسترسی ابرکاربر دارند. \\
\addlinespace
\cmd{/usr/share} & داده‌های اشتراکی (\en{Shared Data}): شامل فایل‌های مستقل از معماری پردازنده (مانند مستندات، فونت‌ها، آیکون‌ها، تم‌ها، فایل‌های راهنما و داده‌های برنامه‌ها). \\
\addlinespace
\cmd{/usr/share/doc} & مستندات نرم‌افزارها (\en{Documentation}): شامل مستندات کامل بسته‌های نصب‌شده، از جمله فایل‌های \en{README}، مجوزها، تاریخچه تغییرات و راهنماهای کاربری. \\
\addlinespace
\cmd{/var} & داده‌های متغیر (\en{Variable Data}): محلی برای فایل‌هایی که در طول زمان تغییر می‌کنند؛ مانند کش برنامه‌ها، صف‌های چاپ، پایگاه‌های داده، ایمیل‌ها و گزارش‌های سیستم. \\
\addlinespace
\cmd{/var/log} & گزارش‌های سیستم (\en{System Logs}): شامل گزارش‌های رویدادها، خطاها و فعالیت‌های سیستم و برنامه‌ها. این مسیر برای عیب‌یابی و نظارت بر سیستم از اهمیت بالایی برخوردار است. \\
\addlinespace
\cmd{\textasciitilde/.config} & پیکربندی کاربر (\en{User Configuration}): یک دایرکتوری پنهان در دایرکتوری خانگی کاربر که تنظیمات اختصاصی برنامه‌های مختلف را برای همان کاربر نگهداری می‌کند (معمولاً بر اساس استاندارد \en{XDG}). \\
\addlinespace
\cmd{\textasciitilde/.local} & داده‌های محلی کاربر (\en{Local User Data}): شامل برنامه‌ها، کتابخانه‌ها و فایل‌های موقتی که فقط برای همان کاربر نصب یا استفاده می‌شوند و نیازی به دسترسی سراسری ندارند. \\

\end{longtable}

\section{مفهوم و کاربرد پیوندهای نمادین (Symbolic Links)}

پیوندهای نمادین (\en{Symbolic Links})، که در ادبیات فنی با نام‌های \en{Soft Link} یا \en{Symlink} نیز شناخته می‌شوند، نوع متمایزی از فایل‌ها در لینوکس هستند که به‌عنوان یک ارجاع یا اشاره‌گر به فایل یا دایرکتوری دیگری عمل می‌کنند. این ویژگی به یک موجودیت واحد در سیستم‌ فایل اجازه می‌دهد تا تحت چندین نام مختلف شناسایی شود؛ قابلیتی که در مدیریت نسخه‌های نرم‌افزاری و ساختاردهی به کتابخانه‌های سیستمی، نقشی راهبردی ایفا می‌کند.

هنگام پیمایش در دایرکتوری‌های سیستمی (نظیر \file{/usr/lib}) با استفاده از دستور \cmd{ls -l}، ممکن است با ورودی‌هایی مواجه شوید که رشته‌ی مجوزهای آن‌ها با حرف \cmd{l} آغاز می‌شود:

\begin{latin}
\begin{lstlisting}
lrwxrwxrwx 1 root root 11 2007-08-11 07:34 libc.so.6 -> libc-2.6.so
\end{lstlisting}
\end{latin}

در خروجی دستور \cmd{ls -l}، درج کاراکتر \cmd{l} در ابتدای ستون مجوزها، نشان‌دهنده‌ی ماهیت پیوندیِ آن فایل است. پیوندهای نمادین در سیستم‌های شبه‌یونیکس، عملکردی مشابه میان‌برها (\en{Shortcuts}) در سایر سیستم‌عامل‌ها دارند، با این تفاوت که در سطح هسته و سیستم‌ فایل به‌صورت یکپارچه پشتیبانی شده و قابلیت‌های مدیریتی بسیار پیشرفته‌تری را برای سازمان‌دهی منابع فراهم می‌کنند.

تصور کنید یک کتابخانه‌ی نرم‌افزاری حیاتی به نام \cmd{foo} دارای نسخه‌های متعددی (مانند \cmd{foo-2.6} و \cmd{foo-2.7}) است. برای مدیران سیستم، درج شماره‌ی نسخه در نام فایل جهت پایش دقیق نسخه‌ی فعال ضرورت دارد. با این حال، تغییر نام فایل اصلی می‌تواند منجر به از کار افتادن تمامی برنامه‌های وابسته شود؛ چرا که به‌روزرسانی تک‌تک این ارجاعات، فرآیندی زمان‌بر و مستعد خطا است.

در چنین سناریویی، پیوندهای نمادین راه‌حلی هوشمندانه ارائه می‌دهند. با ایجاد یک پیوند نمادین با نام عمومی (مانند \cmd{foo}) که به نسخه‌ی واقعی (\cmd{foo-2.6}) اشاره می‌کند، برنامه‌ها همواره به نام ثابت \cmd{foo} ارجاع می‌دهند و سیستم به‌صورت خودکار آن‌ها را به فایل مقصد هدایت می‌کند.

در فرآیند ارتقا به نسخه‌ی جدید (\cmd{foo-2.7})، تنها طی کردن مراحل زیر کفایت می‌کند:

\begin{enumerate}
  \item استقرار فایل نسخه‌ی جدید در مسیر مورد نظر.
  \item حذف پیوند نمادین قدیمی.
  \item ایجاد یک پیوند نمادین جدید با همان نام عمومی (\cmd{foo}) که این بار به نسخه‌ی جدید (\cmd{foo-2.7}) اشاره دارد.
\end{enumerate}

این متدولوژی نه‌تنها فرآیند به‌روزرسانی را تسهیل می‌کند، بلکه امکان همزیستی چندین نسخه را فراهم می‌سازد. در صورت بروز ناپایداری در نسخه‌ی جدید، می‌توان با تغییر مجدد پیوند، سیستم را به‌سرعت به نسخه‌ی پایدار قبلی بازگرداند.

یک نمونه‌ی استاندارد از این سازوکار در دایرکتوری \file{/usr/lib} قابل مشاهده است. در این مسیر، پیوند نمادینی به نام \file{libc.so.6} به فایل اصلی کتابخانه‌ی سیستمی، یعنی \file{libc-2.6.so} اشاره دارد. تمامی برنامه‌هایی که فراخوانی‌های خود را به \file{libc.so.6} ارسال می‌کنند، در لایه‌ی زیرین با \file{libc-2.6.so} تعامل دارند. این مکانیسم، مدیریت وابستگی‌های نرم‌افزاری را در لینوکس بسیار کارآمد و منعطف می‌سازد.

\section{پیوندهای سخت (Hard Links)}

علاوه بر پیوندهای نمادین، مکانیزم دیگری به نام پیوند سخت (\en{Hard Link}) وجود دارد که امکان تخصیص چندین نام به یک فایل واحد را فراهم می‌آورد. با این حال، پیوندهای سخت از منظر فنی عملکردی کاملاً متمایز از پیوندهای نمادین دارند. در حالی که یک پیوند نمادین صرفاً اشاره‌گری به مسیر فایل اصلی است، یک پیوند سخت در واقع نامی ثانویه برای داده‌های فیزیکی یکسان بر روی دیسک (در سطح \en{Inode}) ایجاد می‌کند. در نتیجه، پیوند سخت و فایل اصلی از دیدگاه سیستم‌عامل کاملاً هم‌ارز و غیرقابل تفکیک هستند. در فصول آتی، تفاوت‌های ساختاری و سناریوهای کاربردی هر دو نوع پیوند را به‌طور تفصیلی کالبدشکافی خواهیم کرد.

\section{جمع‌بندی}

با اتمام این فصل، به درک عمیق‌تری از معماری سیستم‌ فایل لینوکس دست یافتیم و با ساختار دایرکتوری‌ها، دسته‌بندی انواع فایل‌ها و ماهیت محتوای آن‌ها آشنا شدیم. آموزه‌ی کلیدی که باید همواره مدنظر داشت، شفافیت و باز بودن ساختار لینوکس است. برخلاف بسیاری از سیستم‌عامل‌های انحصاری (\en{Proprietary})، بخش اعظم فایل‌های راهبردی سیستم، به‌ویژه اسکریپت‌ها و فایل‌های پیکربندی، در قالب متن ساده ذخیره شده‌اند. این رویکرد به کاربران اجازه می‌دهد تا با مطالعه‌ی مستقیم این مستندات، به لایه‌های درونی سیستم نفوذ کرده و منطق عملکردی آن را به‌خوبی درک کنند.

\section{منابع مطالعه تکمیلی}

نسخه کامل استاندارد سلسله‌مراتب سیستم فایل لینوکس (\en{Filesystem Hierarchy Standard}) را می‌توانید در اینجا بیابید:

\begin{latin}
\url{https://refspecs.linuxfoundation.org/fhs.shtml}
\end{latin}

مقاله‌ای پیرامون ساختار دایرکتوری در سیستم‌عامل‌های یونیکس (\en{Unix}) و شبه‌یونیکس (\en{Unix-like}):

\begin{latin}
\url{http://en.wikipedia.org/wiki/Unix_directory_structure}
\end{latin}

شرح مفصلی از قالب متنی اسکی (\en{ASCII}):

\begin{latin}
\url{http://en.wikipedia.org/wiki/ASCII}
\end{latin}